\section*{अध्याय \ref{ch:b2:limb-organogenesis} --- अंगजनन: चतुष्पाद अंग का निर्माण}

\begin{solution}{exo:b2:limb-organogenesis:1}
समीपस्थ-दूरस्थ: सिरे के साथ \omterm{def:b2:limb-organogenesis:bud}{शीर्षस्थ बाह्यचर्मी कटक}, FGF।
अग्र-पश्च: पिछले किनारे पर \omterm{def:b2:limb-organogenesis:bud}{ध्रुवणकारी सक्रियता क्षेत्र},
\omterm{prop:b2:limb-organogenesis:zpa}{Sonic hedgehog}। पृष्ठ-अधर: पृष्ठीय बाह्यचर्म, Wnt7a।
\end{solution}

\begin{solution}{exo:b2:limb-organogenesis:2}
स्टाइलोपोड: प्रगंडिका, ऊर्विका (Hoxa/d 9--10)। ज़्यूगोपोड: अरत्निका और
अंतःप्रकोष्ठिका, अंतर्जंघिका और बहिर्जंघिका (Hox 11)। ऑटोपोड: कलाई और
हाथ, टख़ना और पैर (Hox 13)।
\end{solution}

\begin{solution}{exo:b2:limb-organogenesis:3}
अंग-क्षेत्र निर्दिष्ट (Tbx5); कटक की FGF के नीचे कलिका उभरती है;
ZPA और पृष्ठीय बाह्यचर्म उसे ढाँचा देते हैं; \omterm{prop:b2:limb-organogenesis:aer}{प्रगति क्षेत्र} छोड़ती
कोशिकाएँ किरणों में संघनित होती हैं; उपास्थि विभेदित होती है; संधियाँ
वहाँ बनती हैं जहाँ उपास्थि दबाई जाती है; रक्त वाहिकाएँ भीतर आती हैं और
उपास्थि की जगह अस्थि ले लेती है; अंतरअंगुलीय कोशिकाएँ मरती हैं; और
कायखंडों से पेशी-पूर्वज भीतर प्रवास करते, संलयित होते और कंडराओं से जुड़ जाते हैं।
\end{solution}

\begin{solution}{exo:b2:limb-organogenesis:4}
(क) थोड़ी या बिना प्रगंडिका का ठूँठ; (ख) प्रगंडिका, अरत्निका और
अंतःप्रकोष्ठिका, पर कोई हाथ नहीं; (ग) पूरा अंग --- FGF कटक की जगह ले लेती है।
\end{solution}

\begin{solution}{exo:b2:limb-organogenesis:5}
$x = 40\ln 2 = \qty{28}{\micro m}$, $40\ln 5 = \qty{64}{\micro m}$,
$40\ln 20 = \qty{120}{\micro m}$: यानी 28, 36 और
\qty{56}{\micro m} के क्षेत्र।
\end{solution}

\begin{solution}{exo:b2:limb-organogenesis:6}
(क) 4-3-2-2-3-4, यानी पूरा दर्पण-समुच्चय। (ख) कमज़ोर स्रोत
नीचा शिखर देता है, इसलिए केवल नीची-देहली वाली अंगुली जुड़ती है: 2-2-3-4। (ग)
छोटा अनावरण भी वैसे ही केवल सबसे अग्र अतिरिक्त
अंगुली निर्दिष्ट करता है: 2-2-3-4 --- कोशिकाएँ सांद्रता को काल पर समाकलित करती हैं।
\end{solution}

\begin{solution}{exo:b2:limb-organogenesis:7}
$5000\times 2^{8} = 1.3\times 10^{6}$। $4\times 10^{6}$ तक पहुँचने को
$\log_2 800 = 9.6$ दुगुनेपन चाहिए, इसलिए उस चक्र का औसत लगभग
\qty{10}{h} होना चाहिए, या उस गिनती में वे पेशी-पूर्वज भी होने चाहिए जो
कायखंडों से भीतर प्रवास करते हैं; जबकि कोशिका-मृत्यु उस विसंगति को
बढ़ाती, घटाती नहीं।
\end{solution}

\begin{solution}{exo:b2:limb-organogenesis:8}
चूज़े की अंतरअंगुलीय कोशिकाओं को BMP संकेत मिलता है और वे मर जाती हैं;
जबकि बत्तख़ की अंतरअंगुलीय कोशिकाएँ कोई BMP संदमक अभिव्यक्त करती हैं और
बची रहती हैं, इसलिए वह जाली टिकी रहती है। मनके पर वही संदमक चूज़े की जाली
बचा लेता है: यानी जालीदार पैर। ग्राफ़्ट में भी समय पर मृत्यु का अर्थ है
कि उन कोशिकाओं को हटाए जाने से पहले ही निर्देश मिल चुका था --- यानी वह
निर्णय जल्दी लिया जाता है और बाद में, स्वायत्त रूप से, कार्यान्वित होता है।
\end{solution}

\begin{solution}{exo:b2:limb-organogenesis:9}
Hoxa13 और Hoxd13 के बिना: कोई ऑटोपोड नहीं --- वह अग्र-अंग
कलाई पर समाप्त हो जाता है। Hoxa11 और Hoxd11 के बिना: बहुत छोटी अरत्निका
तथा अंतःप्रकोष्ठिका, और उन पर लगभग सीधे प्रगंडिका से जुड़ा एक हाथ। हर खंड
को अपनी ही Hox जोड़ी चाहिए; वे प्रांत केवल चिह्न नहीं बल्कि
निर्देश हैं।
\end{solution}

\begin{solution}{exo:b2:limb-organogenesis:10}
कटक की FGF ZPA में Shh चालू रखती है; और Shh, किसी मध्यचर्मी
रिले के रास्ते, कटक में FGF चालू रखता है। इसलिए हर केंद्र का संकेत
प्रमाणित करता है कि दूसरा काम कर रहा है, इसलिए कोई अंग बिना ढाँचे के
कभी नहीं बढ़ता और न बिना बढ़े उसे ढाँचा मिलता है, और किसी एक की क्षणिक
हानि दूसरा सुधार देता है। यदि Shh FGF को बनाए न रख सके,
तो वह कटक दिन भर में क्षीण हो जाएगा और उभार जिस भी खंड तक पहुँचा
होगा वहीं रुक जाएगा: यानी एक कटा हुआ अंग। प्रसव का
चक्र उसी घटना से समाप्त होता है जिसे वह चलाता है --- प्रसव उस
उद्दीपन को हटा देता है; जबकि अंग का चक्र विभेदन से समाप्त होता है ---
अंग पूरा होते ही मध्यचर्म रिले करना बंद कर देता है।
\end{solution}

\begin{solution}{exo:b2:limb-organogenesis:11}
$e^{4r} = 8$: $r = \ln 8/4 = \qty{0.52}{d^{-1}}$। वृद्धि: अंगुली की
उपास्थि का लंबा, तेज़ प्रचुरोद्भवन (कोई अंग-संवर्धक जो अंगुलियों में
किसी वृद्धि-जीन की अभिव्यक्ति चढ़ा दे)। मृत्यु: किसी BMP संदमक की
अभिव्यक्ति से अंतरअंगुलीय ऊतक का बचा रहना,
जो वह झिल्ली बनाए रखता है।
\end{solution}

\begin{solution}{exo:b2:limb-organogenesis:12}
संकेत इसलिए दोहराए जाते हैं कि कोई काम करता ग्राही--पारेषण मॉड्यूल
फिर से तैनात करना सस्ता है: ऊतकों के बीच जो बदलता है वह वह
संकेत नहीं बल्कि उन जीनों का समुच्चय है जिन्हें ग्रहण करती कोशिकाएँ चालू
करने में सक्षम हैं, इसलिए वही सांद्रता-परिच्छेद कहीं प्रेरक न्यूरॉन और
कहीं अंगुली 4 का अर्थ रखता है। इसलिए ऐसे किसी जीन का \omterm{def:b2:genome-diversification:mutation}{उत्परिवर्तन}
एक साथ अनेक अंगों को प्रभावित करता है --- Shh के \omterm{def:b2:genome-diversification:mutation}{उत्परिवर्तन}
होलोप्रोसेन्सेफ़ली और अंग-दोष साथ-साथ देते हैं --- जब तक वह किसी
ऊतक-विशिष्ट संवर्धक में न पड़े, और तब वह केवल एक ही अंग बदलता है,
और इन जीनों में अधिकांश विकासात्मक परिवर्तन इसी तरह होता है।
\end{solution}

\begin{solution}{pb:b2:limb-organogenesis:1}
\textbf{1.} $96/12 = 8$ दुगुनेपन: $2000\times 256 = 5.1\times 10^{5}$
कोशिकाएँ।
\textbf{2.} $1000/256 \approx 4$: यानी कोशिकाओं का बड़ा होना और
उपास्थि-आधात्री।
\textbf{3.} लंबाई $\times 10$; अचर त्रिज्या का कोई बेलन आयतन में
$\times 10$ देता; और बाक़ी सौ गुने का अर्थ है कि त्रिज्या भी
कोई दस गुना बढ़ी --- यानी वह कलिका लंबी होते-होते चौड़ी होकर चप्पू बन गई।
\textbf{4.} 3 दिन पर $300/400 = \qty{75}{\%}$; और 7 दिन पर $300/4000 =
\qty{7.5}{\%}$।
\textbf{5.} केवल सिरा ही उस संकेत के नीचे है; और जैसे-जैसे सिरा आगे
बढ़ता है, पीछे छूटी कोशिकाएँ उसी क्रम में उस क्षेत्र से निकलती हैं जिस
क्रम में वे बनीं, इसलिए समीपस्थ संरचनाएँ पहले और दूरस्थ अंत में
निर्दिष्ट होती हैं।
\textbf{6.} ज़्यूगोपोड (अरत्निका और अंतःप्रकोष्ठिका), जो 3.5 और
4 दिन के बीच निर्दिष्ट हुआ।
\textbf{7.} स्टाइलोपोड 3.5 दिन तक, ज़्यूगोपोड 4.0 दिन तक, ऑटोपोड 4.5
दिन तक (अंगुलियों के सिरे बाद में)।
\textbf{8.} आधा दिन: यानी एक कोशिका-चक्र।
\textbf{9.} \omterm{prop:b2:limb-organogenesis:aer}{प्रगति क्षेत्र} में बिताया हर चक्र उस Hox
कूट को आगे बढ़ाता है (9--10, फिर 11, फिर 13); और एक, दो या
तीन चक्र बाद निकलती कोशिकाएँ स्टाइलोपोड, ज़्यूगोपोड या ऑटोपोड का कूट ढोती हैं।
\textbf{10.} पूरा पंख: कटक जो देता था वह FGF ही थी।
\textbf{11.} पृष्ठीय सतह से एक दूसरा उभार: यानी दोहराई गई
दूरस्थ संरचनाएँ, एक अतिरिक्त अंग।
\textbf{12.} मध्यचर्म कटक तक Shh रिले करना बंद कर देता है, Shh की
अभिव्यक्ति समाप्त हो जाती है, कटक क्षीण हो जाता है, और कोशिकाएँ विभाजित
होने के बजाय विभेदित हो जाती हैं: यानी वह चक्र टूट जाता है जो वृद्धि को बनाए रखता था।
\textbf{13.} $k = \ln 2/1740 = \qty{4.0e-4}{s^{-1}}$; $\lambda =
\sqrt{10^{-12}/4\times 10^{-4}} = \qty{50}{\micro m}$।
\textbf{14.} $x = 50\ln(1/0.6) = 26$, $50\ln 4 = 69$, $50\ln 12.5 =
\qty{126}{\micro m}$: यानी 26, 43 और \qty{57}{\micro m} के क्षेत्र।
\textbf{15.} $600 - 126 = \qty{474}{\micro m}$।
\textbf{16.} हर सीमा $50\ln 2 = \qty{35}{\micro
m}$ बाहर खिसक जाती है: 61, 104, 161। तीनों चौड़ाइयाँ 61, 43 और \qty{57}{\micro m} हैं:
यानी केवल अंगुली-4 का क्षेत्र बढ़ता है, बाक़ी केवल खिसक जाते हैं।
वह प्रतिरूप कोई अंगुली पाने के बजाय अग्र की ओर खिसक जाता है, और
अंगुली-रहित अग्र क्षेत्र 474 से घटकर \qty{439}{\micro m} रह जाता है।
\textbf{17.} हर समुच्चय को \qty{126}{\micro m} चाहिए: यानी कम से कम
\qty{252}{\micro m}; और \qty{600}{\micro m} की कलिका में जगह है, जिसके
बीच में \qty{348}{\micro m} अंगुली-रहित ऊतक बचता है।
\textbf{18.} \qty{200}{\micro m} पर दोनों प्रवणताएँ अतिव्यापित होकर
जुड़ जाती हैं: बीच की कोशिकाएँ अंगुली 3 लायक़ देख लेती हैं और अंगुली-2 के
क्षेत्र लुप्त हो जाते हैं --- यानी 4-3-4 या 4-3-3-4, किसी पूरे
दर्पण-समुच्चय से कम अंगुलियाँ।
\textbf{19.} चौथाई कोशिकाएँ $c_0/4$ देती हैं: कोई कोशिका
अंगुली 4 या 3 की देहली तक नहीं पहुँचती (किनारे पर ही $0.25 c_0$), और
अंगुली-2 की देहली $50\ln(0.25/0.08) =
\qty{57}{\micro m}$ तक पार होती है: यानी एक अतिरिक्त अंगुली 2।
\textbf{20.} $3\times 300\times 300\times 200 = 5.4\times 10^{7}$
\unit{\micro m^3}: यानी $5.4\times 10^{4}$ कोशिकाएँ; और \qty{86400}{s} में
प्रति सेकंड $0.6$ कोशिकाएँ।
\textbf{21.} बत्तख़ की जाली में BMP का मनका: वह जाली मर जाती है और
अंगुलियाँ अलग हो जाती हैं; चूज़े की जाली में संदमक: वह जाली बच जाती है, यानी जालीदार
पैर।
\textbf{22.} $e^{0.5\times 4} = e^{2} = 7.4$। दूसरा परिवर्तन
किसी BMP संदमक से अंतरअंगुलीय झिल्ली का बचा रहना है।
\textbf{23.} Hox 13 की वह पश्च प्रावस्था जो ऑटोपोड निर्दिष्ट करती है,
और उस कटक का बना रहना: मछली की तह उसके बजाय पख-किरणें
बनाती है। \emph{Tiktaalik} के पख में कलाई जैसी अस्थियाँ हैं --- यानी
अंग बनता हुआ एक पख।
\textbf{24.} $5.4\times 10^{4}/1000 = 54$: $\log_2 54 = 5.8$ दुगुनेपन,
यानी लगभग \qty{69}{h}, तीन दिन --- जिन्हें हटाने में एक दिन लगता है।
\textbf{25.} आठ दुगुनेपन; अंगुली-सीमाएँ 26, 69 और
\qty{126}{\micro m} पर; और किसी दर्पण-द्विगुणन को कम से कम
\qty{252}{\micro m} चाहिए।
\end{solution}
